\part{Répertoires et variables d'environnement}
\label{part:arbo}
\section{Répertoires}
\label{sec:repert}
   Tous les fichiers et programmes qui constituent AIDA se regroupent de
manière arborescente sous un unique répertoire racine. Il n'y a pas 
de contrainte sur le choix de l'emplacement de ce répertoire, qui peut 
varier d'une machine à l'autre.

 Dans ce répertoire-racine, se trouvent les sous-répertoires suivants :
\begin{description}
  \item[src] contient les sources des programmes qui constituent AIDA. Il 
   est divisés en sous-répertoires, un par langage (script, 
   python3, etc), plus un pour les bibliothèques générales écrites en C
   (\texttt{lib}) et un contenant les codes sources de la documentation
   AIDA (\texttt{doc}).

  \item[bin] contient tous les exécutables d'AIDA. Pour l'utilisation
   confortable d'AIDA, il est recommandé de rajouter ce répertoire
   dans sa variable PATH (voir \path{~/.profile.aida}) :

  \item[appli] contient le matériel de traitements des campagnes et des
    acquistions de l'utilisateur (scripts, programmes, fichires de
    configuration...). \\
    \textbf{Ce répertoire n'est pas géré par le gestionnaire de source de AIDA.}

  \item[meta] contient tous les fichiers de méta-données. Il contient au moins
   deux sous-répertoires utilisés par AIDA :
      \begin{description}
       \item[INF] : les meta-données des acquisitions sous forme de fichiers
                    <acqid>.inf.
       \item[campagnes] : les meta-données des campagnes sous forme de fichiers
		                <campagne>.camp.
      \end{description}
   L'utilisateur peut y ajouter les sous-répertoires utilisés par son matériel
   de traitement (par exemple, \texttt{crvol} pour les comptes-rendus de vol).
   Si le réseau local le permet, ce peut être un montage NFS ou un
   partage SMB.

  \item[local] contient les fichiers spécifiques à la configuration 
   informatique de la machine sur laquelle s'execute AIDA.

  \item[config] contient un certain nombre de fichiers 
   de configuration d'AIDA (liste des familles de paramètres, les définitions
   de graphiques).
   Ce répertoire peut être un montage NFS ou un partage SMB.

  \item[profils] contient un sous-répertoire spécifique pour chaque profil
   d'utilisation de AIDA (ex: 4M, EXEMPLES, SAFIRE, TRAMM).
   Chaque sous-répertoires contient au moins un fichier \texttt{plate-formes}
   (liste de plate-forme connues pour ce profil d'utilisation de AIDA).

  \item[share] contient des fichiers partageables (cartes, logos...).

  \item[data] contient toutes les données. Pour des raisons pratiques,
   il s'agit souvent d'un lien vers un gros disque dédié aux données.
   Si le réseau local le permet, ce peut être un montage NFS ou un
   partage SMB.
   L'organisation des données est décrite un peu plus loin 
   (paragraphe \ref{sec:arbodata}).

  \item[doc] contient la documentation d'AIDA (donc, en particulier, ce
   fichier). Attention, ce n'est pas la même chose que la documentation
   sur les vols ou les traitements, produite par AIDA.

  \item[include] : contient tous les fichiers include (\path{*.h} du C, ...).

  \item[lib] contient toutes les bibliothèques AIDA (fichiers \path{*.a}, 
   \path{*.sl}, \path{*.ln}, \path{*.so} mais aussi sous le répertoire
   \path{pythonX.Y}, les librairies python de AIDA).

  \item[log] pour les programmes qui veulent laisser une trace de ce qui
   s'est passé pendant leur éxécution.

  \item[extern] accueille le matériel tierces utilisé par AIDA, lorsqu'elles
   ne sont pas fournies par le système (par exemple \texttt{cextract}, ...)
\end{description}
 
\section{Variables d'environnement}
   L'emplacement du répertoire racine pouvant varier, on utilise une 
variable d'environnement, \textsf{AIDA\_HOME}, pour en 
désigner l'emplacement. Tous les programmes prennent en compte cette
variable, il est donc indispensable qu'elle soit établie dans 
l'environnement de chaque utilisateur.

   A l'installation de AIDA, les variables d'environnement suivantes sont
déclarées (voir \path{~/.profile.aida}) :
\begin{description}
  \item[\textsf{AIDA\_HOME}] désigne l'emplacement du répertoire racine de AIDA.
  \item[\textsf{AIDA\_PROFIL}] désigne le profil d'utilisateur AIDA.
  \item[\textsf{AIDA\_META}] désigne l'emplacement du répertoire des
   méta-données.
  \item[\textsf{AIDA\_DATA}] désigne l'emplacement du répertoire des données
\end{description}

\medbreak
   D'autre variables peuvent désigner chacun des sous-répertoires
décrits au paragraphe \ref{sec:repert}. Elle sont utiles lorsqu'un
utilisateur ne souhaite pas utiliser l'arborescence commune mise en place,
par exemple pour accéder à des données situées ailleurs que dans
\texttt{\$AIDA\_HOME\//data}, ou à des versions spécifiques
de programmes situés ailleurs que dans \texttt{\$AIDA\_HOME\//bin}.\\
La règle pour construire le nom de ces variables est simple : elles sont
baptisées \textsf{AIDA\_\textit{xxx}}, où \textit{xxx} est le nom en majuscule
du sous-répertoire qu'elle désigne. \\
Ainsi :
\begin{itemize}
  \item[\texttt{AIDA\_BIN}] : sert à accéder à des programmes situé ailleurs
    que dans \texttt{\$AIDA\_HOME\//bin},
  \item[\texttt{AIDA\_LIB}] : sert à utiliser des librairies situées ailleurs
    que dans \texttt{\$AIDA\_HOME\//lib},
  \item[\texttt{AIDA\_DATA}] : sert à travailler sur des données situées
   ailleurs que dans \texttt{\$AIDA\_HOME\//data}.
\end{itemize}

\section{Identification de chaque période de mesure}
\label{sec:AcqId}
   On appelle \textit{période de mesure} le temps pendant lequel
un système d'acquisition a fonctionné. Pour les systèmes qui
fonctionnent par période, la définition est assez simple (un 
lâcher de ballon-sonde, un vol d'un avion). Pour les systèmes qui
fonctionnent en continu, la période de mesure est en général
définie comme une journée, mais rien n'empèche  de choisir une
durée plus courte ou plus longue (le mois, pas exemple).

\medbreak
On distingue 5 types de période de mesure :

\textbf{- Journee} (\textit{anciennement Acquisition})

\textbf{- Vol}

\textbf{- Lacher}

\textbf{- Lacher\_en\_vol}

\textbf{- Drone}


\medbreak
   Pour des raisons pratiques (nommage de fichiers, choix d'une
période parmis les autres, ..), il est important de pouvoir désigner
chaque période de mesure par un nom unique, appelé \textbf{\texttt{acqid}}.

\medbreak
Un \textbf{\texttt{acqid}} est une chaîne de 8 caractères, de la forme : \\
\begin{center}
{\fbox{ \texttt{ppAAnnnN}}}
\end{center}

\medbreak
\begin{itemize}
  \item Les 2 premiers (\texttt{pp}) sont alpha-numériques, et désignent la 
   plate-forme de mesure. La correspondance entre cette abbréviation et
   le nom usuel, la description, etc se fait par le fichier
   {\texttt plate-formes} décrit au paragraphe \ref{sec:platef}, page 
   \pageref{sec:platef}.

  \medbreak
  \item Les 2 suivants (\texttt{AA}) sont numériques, et désignent l'année.
   Les valeurs supérieures à 70 correspondent aux années 1970 à 1999,
   et les autres aux années 2000 à 2069.

  \medbreak
  \item Les 4 suivants (\texttt{nnnN}) désignent le
   numéro d'acquisition de manière différente selon le type de période de
   mesure :

   \begin{itemize}

     \item[\textbf{Journee}] : \texttt{nnnN} est le quantième du jour dans
                               l'année (numérique)

     \item[\textbf{Vol}] : \texttt{nnnN} est le quantième du vol dans l'année
                           (numérique, sur 4 chiffres commençant par 0)

     \item[\textbf{Lacher}] : \texttt{nnnN} est le nombre d'heures depuis le 01/01
                              de l'année (numérique)

     \item[\textbf{Lacher\_en\_vol}] : \texttt{nnn} est le numéro du vol (numérique),
                                       \texttt{N} est le numéro du lâcher
                                       (alphanumérique [0-9] puis [a-z] puis [A-Z]
                                       soit 62 numéros de lâcher)

     \item[\textbf{Drone}] : \texttt{nnn} est le quantième du jour dans l'année
                             (numérique), \texttt{N} est le numéro du vol (idem
                             numéro de lâcher pour \texttt{Lacher\_en\_vol})
   \end{itemize}
\end{itemize}

\textbf{\textit{Note : avant 2008, la partie \texttt{nnnN} des acqids de type
vol n'était constitué que de 2 chiffres.}}


\section{Détail des arborescences principales}

Certains des sous-répertoires présentés ci-dessus contiennent
une information volumineuse et/ou complexe : ils sont en fait le sommet
d'une arborescence. Nous décrivons ici les plus importantes.

\subsection{L'arborescence des données} 
\label{sec:arbodata}

    Dans le cadre d'un fonctionnement normal, l'utilisateur 
n'a pas besoin de connaître les détails du stockage
des données.

    Les programmes (et les bibliothèques sous-jacentes) sont conçues 
pour qu'il suffise à l'utilisateur de spécifier 
l'identification d'une période de mesure (par exempl
 l'\textbf{\texttt{acqid}} \texttt{Tf210356})
pour que les programmes d'AIDA sachent accéder aux données.

   Pour que cela fonctionne, il faut évidement spécifier une 
organisation de stockage, qui doit être respectée par tous les 
programmes.  Actuellement, cette organisation du stockage repose
sur le système de fichiers et de répertoires arborescents 
fourni par le système d'exploitation. C'est une arborescence 
à trois niveaux : le \textbf{premier correspond à la plate-forme,
le second à l'année, et le troisième à la période de mesure}.

    Cette organisation permet d'avoir une architecture complètement
ouverte, en permettant au plus grand nombre de logiciels
et de langages des programmation de travailler sur les données d'AIDA.

    L'emplacement d'un fichier se déduit automatiquement
du nom unique donné à chaque période de mesure :

    Le répertoire \texttt{\$AIDA\_HOME\//data} est le
point de départ de toute l'arborescence des données. Il contient un
répertoire par plate-forme de mesure. Le nom de ces répertoires
est exactement le même que celui inscrit dans le fichier
\texttt{plate-forme} (décrit au paragraphe \ref{sec:platef}).
A son tour, chacun de ces répertoires contient un répertoire par 
année, dont le nom est écrit en quatre chiffres, et qui contient
un répertoire par période de mesure. Enfin, dans ce répertoire
«feuille», il y a un fichier par
paramètre, qui contient la série complète des données (format
\textit{vecteur} décrit en détail dans la partie suivante).

 En reprennant l'exemple de  l'acquisition \texttt{Tf210356}, le 
répertoire sera : \\
\texttt{\$AIDA\_HOME\//data\//toulouse-flux\//2021\//0356}.

\subsection{L'arborescence des applications} 
    Pour chaque campagne de mesure, il est nécessaire
d'avoir des fichiers particularisés : des \textit{scripts}
de traitements, des scripts de production des fichiers
ASCII, ....

     Ces fichiers correspondent à des \textbf{applications}
utilisant AIDA, et sont regroupés dans \texttt{\$AIDA\_HOME\//appli},
plus précisément dans :
\texttt{\$AIDA\_HOME\//appli\//script\//\$nom\_campagne}.


